\begin{helper}[Integrability of the structure weight along a segment]
Let $n\ge0$ and $p>1$.  Then for every $\varepsilon\ge0$ and all
$a,b\in\R^n$ the function
\[
t\ \longmapsto\ \bigl(\varepsilon+|ta+(1-t)b|^2\bigr)^{\frac p2-1}
\]
is interval-integrable with respect to Lebesgue measure on $[0,1]$.
All powers are real powers of nonnegative reals, with the convention
$0^{s}=0$ for $s\ne0$ and $0^{0}=1$; in the single degenerate configuration
$\varepsilon=0$, $a=b=0$, $p<2$ the integrand is therefore the constant $0$
rather than $+\infty$, and the statement is read with that convention.
\end{helper}
\begin{proof}
Write $s=\frac p2-1$, so that $2s=p-2$ and the hypothesis $p>1$ says exactly
$2s>-1$.  Put
\[
\varphi(t)=ta+(1-t)b=b+t(a-b),\qquad
W(t)=\bigl(\varepsilon+|\varphi(t)|^2\bigr)^{s}\quad(t\in\R).
\]

\emph{Measurability.}  The map $t\mapsto\varepsilon+|\varphi(t)|^2$ is a real
quadratic polynomial in $t$, hence continuous; it takes values in
$[0,\infty)$ because $\varepsilon\ge0$.  The real power $u\mapsto u^{s}$ is
Borel measurable on $\R$.  Hence $W$ is measurable, and in every case below it
remains to produce an integrable majorant on $[0,1]$.

\emph{Case 1: $p\ge2$, i.e.\ $s\ge0$.}  If $s=0$ then $W\equiv1$ by the
convention $u^{0}=1$, which is continuous.  If $s>0$ then $u\mapsto u^{s}$ is
continuous on $[0,\infty)$, so $W$ is continuous on $[0,1]$ as a composition of
continuous maps.  A continuous function on a compact interval is
interval-integrable, which settles this case.

\emph{Case 2: $1<p<2$ and $\varepsilon>0$, i.e.\ $s<0$ with a positive floor.}
For every $t$ we have $\varepsilon+|\varphi(t)|^2\ge\varepsilon>0$, so
$t\mapsto\varepsilon+|\varphi(t)|^2$ takes values in $[\varepsilon,\infty)$,
on which $u\mapsto u^{s}$ is continuous.  Hence $W$ is continuous on $[0,1]$
and therefore interval-integrable.  (Quantitatively $0\le W\le\varepsilon^{s}$,
since $s<0$ and the base is $\ge\varepsilon$.)

\emph{Case 3: $1<p<2$ and $\varepsilon=0$, i.e.\ $s<0$ and no floor.}  This is
the degenerate corner, and $W(t)=|\varphi(t)|^{2s}=|\varphi(t)|^{p-2}$.

\emph{Case 3a: $a=b$.}  Then $\varphi\equiv b$ and $W$ is the constant
$|b|^{p-2}$ (which is the constant $0$ if $b=0$, by the stated convention).  A
constant is interval-integrable.

\emph{Case 3b: $a\ne b$.}  Set
\[
A=|a-b|^2>0,\qquad B=(b,a-b),\qquad C=|b|^2 ,
\]
so that, expanding the inner product,
\[
|\varphi(t)|^2=|b+t(a-b)|^2=At^2+2Bt+C .
\]
Cauchy--Schwarz (Preamble) gives $B^2=(b,a-b)^2\le|b|^2|a-b|^2=AC$.  Let
$t_0=-B/A$.  Completing the square,
\[
At^2+2Bt+C=A\Bigl(t+\frac BA\Bigr)^2+\Bigl(C-\frac{B^2}A\Bigr)
=A(t-t_0)^2+\Bigl(C-\frac{B^2}A\Bigr),
\]
and $C-\frac{B^2}A\ge0$ by the Cauchy--Schwarz inequality just recorded, since
$A>0$.  Therefore
\[
A(t-t_0)^2\ \le\ |\varphi(t)|^2\qquad\text{for every }t\in\R .
\]
Fix $t\ne t_0$.  Then $0<A(t-t_0)^2\le|\varphi(t)|^2$, and since $s<0$ the map
$u\mapsto u^{s}$ is decreasing on $(0,\infty)$, so
\[
W(t)=\bigl(|\varphi(t)|^2\bigr)^{s}\ \le\ \bigl(A(t-t_0)^2\bigr)^{s}
=A^{s}\,|t-t_0|^{2s}=A^{s}\,|t-t_0|^{p-2},
\]
where we used that both bases are positive, so the real-power identities
$(uv)^{s}=u^{s}v^{s}$ and $(u^2)^{s}=|u|^{2s}$ apply.  Also $W\ge0$
everywhere, because a real power of a nonnegative base is nonnegative.  Thus
\[
0\ \le\ W(t)\ \le\ A^{s}\,|t-t_0|^{p-2}
\qquad\text{for all }t\in[0,1]\setminus\{t_0\},
\]
and $\{t_0\}$ is a Lebesgue null set, so the bound holds almost everywhere on
$[0,1]$.

It remains to see that $t\mapsto|t-t_0|^{p-2}$ is interval-integrable on
$[0,1]$.  Since $p>1$ we have $p-2>-1$, so by the Preamble item on
interval integrability of real powers, $u\mapsto u^{p-2}$ is
interval-integrable on every interval; and Lebesgue measure on $\R$ is
invariant under the translation $u=t-t_0$ and under the reflection
$u\mapsto-u$, both of which preserve interval integrability.  We distinguish
the position of $t_0$:

\begin{itemize}
\item If $t_0\le0$, then $t-t_0\ge0$ on $[0,1]$, so
  $|t-t_0|^{p-2}=(t-t_0)^{p-2}$, and the translation $u=t-t_0$ turns the
  integral into $\int_{-t_0}^{1-t_0}u^{p-2}\,du$, finite by the Preamble item.
\item If $t_0\ge1$, then $t_0-t\ge0$ on $[0,1]$, so
  $|t-t_0|^{p-2}=(t_0-t)^{p-2}$, and the reflected translation $u=t_0-t$ turns
  the integral into $\int_{t_0-1}^{t_0}u^{p-2}\,du$, again finite.
\item If $0<t_0<1$, split $[0,1]=[0,t_0]\cup[t_0,1]$.  On $[t_0,1]$ the first
  bullet's computation applies and on $[0,t_0]$ the second bullet's does, and
  interval integrability on two abutting intervals gives interval
  integrability on their union.
\end{itemize}

So $t\mapsto A^{s}|t-t_0|^{p-2}$ is a nonnegative integrable majorant of the
measurable function $W$ on $[0,1]$, whence $W$ is interval-integrable on
$[0,1]$.  This exhausts all cases.  The hypothesis $p>1$ was used only in
Case 3b, to get $p-2>-1$, and it is sharp there: if $\varepsilon=0$, $p\le1$
and the line through $a$ and $b$ passes through the origin at a parameter
$t_0\in[0,1]$ -- i.e.\ $C=B^2/A$ in the notation above, so that
$|\varphi(t)|^2=A(t-t_0)^2$ exactly -- then
$W(t)=A^{s}|t-t_0|^{p-2}$ with $p-2\le-1$, whose integral over $[0,1]$
diverges.
\end{proof}
